\section{线性回归从投影到概率模型}
\label{sec:13-Machine-Learning/02-Linear-Models/01}

\subsection{同小区同户型的两套房，模型给了同一个数}

租房平台用面积、楼龄和到地铁的步行分钟数估计月租金，五百条成交记录拟合下来，训练 MSE 是 \(2840\)，均方根误差约 53 元。汇总数字看着可用。翻明细时冒出一对刺眼的样本：同一小区、同一户型、同一楼层，面积都是 89 平方米，楼龄都是 5 年，步行都是 8 分钟，成交租金却是 2800 元和 4000 元，模型对两套都预测 3500 元上下。产品经理问，这算不算模型出错。

更麻烦的是接下来那一步。团队补上楼层、朝向、装修等级和小区均价，训练 MSE 降到 2100，看着是改善；面积的系数却从 \(+15.3\) 掉到 \(+2.1\)，小区均价的系数是 \(+0.72\)。每多一平方米只值两块钱，这个数字没有任何从业者会认。换个特征加入顺序，先放小区均价再放面积，面积的系数还可能翻成负数。预测值几乎没动，可那个要拿去向业务解释的数字翻了脸。

这两件事出自同一个几何结构。本节先把这个结构讲清楚，再说明概率假设是在哪一步加进来的，以及加它进来要认下哪些前提。

\subsection{最小二乘先是几何，概率是后来才加的}

求最小二乘拟合不需要任何分布假设。它只做一件事：
\[
\widehat\beta\in\argmin_\beta\norm{y-X\beta}_2^2,
\]
其中 \(X\in\R^{n\times p}\) 的每行是一条样本、每列是一项特征，\(y\in\R^n\) 是响应向量，截距靠给 \(X\) 添一列全 1 向量容纳。

所有可能的拟合值 \(X\beta\) 组成 \(X\) 的列空间，它是 \(\R^n\) 里一个维数至多为 \(p\) 的子空间。最小二乘在这个子空间里找离 \(y\) 欧氏距离最近的点，也就是把 \(y\) 正交投影下去。残差 \(r=y-X\widehat\beta\) 必须与列空间中每个方向垂直，否则沿着残差还有分量的那个方向挪一点投影点，距离就能再短一截，与``最近''矛盾。写成矩阵语言即 \(X^{\trans}r=0\)，代入展开得到
\[
X^{\trans}(y-X\widehat\beta)=0
\quad\Longrightarrow\quad
X^{\trans}X\widehat\beta=X^{\trans}y.
\]

这套几何在线性代数里已经建立过，\hyperref[sec:03-Linear-Algebra/03-Orthogonality/02]{``投影与最近点：从一条直线走向一般子空间''}讲的是同一件事的抽象版本，\hyperref[sec:03-Linear-Algebra/03-Orthogonality/03]{``最小二乘：为无解方程寻找最佳近似''}给出了正规方程的推导。往更远处推一步，实分析把它搬进函数空间：\hyperref[sec:12-Real-Analysis-Support/07-Lp-Spaces-and-Function-Geometry/04]{``条件期望与线性回归是两种投影''}说明条件期望 \(\E(Y\given X)\) 是 \(Y\) 在全体 \(X\) 可测函数上的投影，而线性回归只投到其中的线性子空间上。两者都在 \(L^2\) 里取最近点，区别只在允许使用的候选集合有多大。

\begin{center}
\begin{tikzpicture}[>=stealth]
  \fill[black!6] (-3,-0.9) -- (2.4,-1.7) -- (3,0.9) -- (-2.4,1.7) -- cycle;
  \draw[black!45] (-3,-0.9) -- (2.4,-1.7) -- (3,0.9) -- (-2.4,1.7) -- cycle;
  \node[black!55] at (-1.75,-1.52) {\footnotesize \(X\) 的列空间};
  \fill (0,0) circle (1.6pt);
  \node[below left] at (0,0) {\footnotesize \(0\)};
  \draw[->,thick] (0,0) -- (2.05,-0.62);
  \draw[->,thick] (0,0) -- (2.35,-0.42);
  \node[right] at (2.4,-0.36) {\footnotesize \(x_2\)};
  \node[below] at (2.0,-0.72) {\footnotesize \(x_1\)};
  \fill (0.85,0.15) circle (1.6pt);
  \node[anchor=east] at (0.74,0.17) {\footnotesize \(X\widehat\beta\)};
  \draw[->,very thick] (0,0) -- (0.85,2.55);
  \node[above] at (0.85,2.55) {\footnotesize \(y\)};
  \draw[dashed,thick] (0.85,2.55) -- (0.85,0.15);
  \node[right] at (0.92,1.45) {\footnotesize \(r=y-X\widehat\beta\)};
  \draw (0.85,0.37) -- (1.09,0.34) -- (1.07,0.11);
\end{tikzpicture}
\end{center}

\emph{投影点唯一；但当 \(x_1\) 与 \(x_2\) 几乎同向时，把这个点写成两者组合的方式几乎不唯一。}

三个点的手算能让正交条件落在具体数字上。取中心化后的特征 \(x=(-1,0,1)^{\trans}\)，响应 \(y=(1,2,2)^{\trans}\)，拟合 \(\widehat y=\beta_0+\beta_1x\)。截距是响应均值 \(\widehat\beta_0=5/3\)，斜率
\[
\widehat\beta_1=\frac{x^{\trans}y}{x^{\trans}x}
=\frac{(-1)\times 1+0\times 2+1\times 2}{1+0+1}=\frac12 .
\]
拟合值为 \((7/6,\ 5/3,\ 13/6)^{\trans}\)，残差为 \((-1/6,\ 1/3,\ -1/6)^{\trans}\)。残差之和为零，说明它与全 1 列正交；它与 \(x\) 的内积 \((-1/6)(-1)+0+(-1/6)(1)=0\)，说明它与特征列正交。这两个等式是一阶最优条件在这份数据上的逐条兑现。

\subsection{正规方程描述答案，但不承担求解}

\(X^{\trans}X\widehat\beta=X^{\trans}y\) 精确刻画了解满足的关系，把它当成计算指令则是另一回事。形成 \(X^{\trans}X\) 会把条件数平方：若 \(\kappa(X)=10^4\)，则 \(\kappa(X^{\trans}X)\approx 10^8\)，双精度下有效数字从约 16 位掉到约 8 位。教科书写 \(\widehat\beta=(X^{\trans}X)^{-1}X^{\trans}y\) 是为了看清结构，数值库不这么算。

实现走 QR：把 \(X=QR\) 代入残差范数，\(\norm{y-X\beta}_2=\norm{Q^{\trans}y-R\beta}_2\)（正交矩阵保范数），于是只需回代解上三角系统 \(R\widehat\beta=Q^{\trans}y\)，全程只跟 \(\kappa(X)\) 打交道。特征几乎线性相关、\(X\) 近乎秩亏时，QR 的列主元版本或者 SVD 伪逆更稳，算起来也更慢。这条路线的细节在 \hyperref[sec:08-Numerical-Analysis/05-Numerical-Linear-Algebra/02]{``QR 与最小二乘的稳定解法''}与 \hyperref[sec:08-Numerical-Analysis/05-Numerical-Linear-Algebra/06]{``SVD、低秩与数值秩''}里；条件数为什么衡量的是问题本身而非算法，见 \hyperref[sec:08-Numerical-Analysis/01-Floating-Point-and-Error/03]{``条件数衡量问题本身的敏感性''}。

值得留意的是，条件数变差与系数不稳定是同一件事的两副面孔。\(X^{\trans}X\) 的小特征值对应数据中信息稀薄的方向，而 \(\widehat\beta\) 的协方差矩阵正是 \(\sigma^2(X^{\trans}X)^{-1}\)，小特征值在这里被倒过来放大成巨大的方差。换更好的算法能救回浮点精度，救不回数据本身没有提供的区分能力。

\subsection{三层假设各自买到一件不同的东西}

投影已经给出了 \(\widehat\beta\)、拟合值和残差，此时还没有出现任何概率。要回答``这个系数值得多大程度的信任''，才需要往上加假设，而且是分层加的。

第一层假设误差的条件均值为零：\(Y=X\beta+\varepsilon\)，\(\E(\varepsilon\given X)=0\)。这一层买到的是无偏性。骨架很短：\(\widehat\beta=(X^{\trans}X)^{-1}X^{\trans}Y=\beta+(X^{\trans}X)^{-1}X^{\trans}\varepsilon\)，取条件期望，第二项归零。条件对账：这里用到 \(X^{\trans}X\) 可逆（列满秩），用到 \(\E(\varepsilon\given X)=0\) 而非仅仅 \(\E\varepsilon=0\)——后者允许误差随 \(X\) 系统性地偏移，那样第二项就不会消失。还要看清``无偏''说的是重复抽样下的平均行为，它不保证手头这一次的估计接近 \(\beta\)。

第二层再加同方差与不相关：\(\Var(\varepsilon\given X)=\sigma^2I\)。这一层买到方差公式 \(\Var(\widehat\beta\given X)=\sigma^2(X^{\trans}X)^{-1}\)，以及 Gauss--Markov 定理——在所有对 \(y\) 线性且无偏的估计量里，最小二乘的方差最小。骨架是取任一线性无偏估计 \(\tilde\beta=Cy\)，令 \(D=C-(X^{\trans}X)^{-1}X^{\trans}\)，无偏性迫使 \(DX=0\)，于是 \(\Var(\tilde\beta)=\sigma^2(X^{\trans}X)^{-1}+\sigma^2DD^{\trans}\)，多出来的那项半正定。条件对账：结论被``线性''和``无偏''两个词框住了，走出这个类别，有偏而方差更小的估计量可以给出更低的均方误差，下一节的 Ridge 就是。同方差本身是可检验的判断，高价房与低价房的租金波动幅度是否相同，看残差图；同一房东名下的多套房源违反不相关。

第三层假设正态：\(\varepsilon\given X\sim\mathcal N(0,\sigma^2I)\)。这一层买到有限样本的精确分布——最小二乘同时是最大似然估计，\(\widehat\beta\) 服从多元正态，\(t\) 检验与 \(F\) 检验不必诉诸渐近近似。它也是最强、最容易被真实数据违反的一条。丢掉正态性，前两层的结论一条不少，只是小样本的 \(p\) 值需要改用渐近理论或者重采样来取。

\subsection{共线时投影稳定，系数不稳定}

回到开头的系数翻脸。面积与小区均价高度相关，大户型多出现在高端小区。在图里，这对应 \(x_1\) 与 \(x_2\) 两个箭头几乎同向：它们张成的仍是同一张平面，\(y\) 到平面的投影点分毫不差，但把投影点写成 \(\beta_1x_1+\beta_2x_2\) 的方式几乎不唯一——沿着 \(x_1\approx cx_2\) 这个近似关系，可以给 \(\beta_1\) 加上一个量、同时给 \(\beta_2\) 减去相应的量，拟合值几乎不动。数据没有能力在这两条路之间做出选择，训练集的偶然波动就替它选了。

用可识别性来表述会更准确：``面积每增加一平方米，租金增加 \(\beta_1\) 元''这句话要成立，前提是不存在另一组系数产生同样的拟合。共线恰好提供了这种替代自由。诊断上有两个信号，一个是方差膨胀因子 \(\mathrm{VIF}_j=1/(1-R_j^2)\)，其中 \(R_j^2\) 是把第 \(j\) 列对其余列回归的决定系数；另一个是 \(X\) 的条件数。两者都超过经验阈值时，别急着解释单个系数，先问这个系数是否被数据确定。

预测和解释在这里分了岔。共线不伤害同分布下的预测——投影点是稳的；它伤害的是每个变量各自贡献多少这类问句。若部署时特征之间的相关结构发生变化（新城区的大户型不再集中在高端小区），那条被数据偶然选中的系数分配也会连预测一起带偏。

\subsection{系数不是因果效应，哪怕它显著}

最容易被搬到 PPT 上的误读是：面积系数 \(+15.3\) 且 \(p<0.01\)，所以``把房子扩建一平方米，租金会涨 15.3 元''。回归系数说的是在同一份观测数据里、控制住已放入模型的那些列之后，面积每高一单位对应的租金条件均值差异。它描述的是已经发生的数据里变量怎样共同变化，不描述你去动其中一个变量会发生什么。

差别来自三处。其一是未观测混杂：装修等级没进模型，而它同时影响面积（大户型倾向精装）和租金，面积系数就吞掉了装修的一部分效应。其二是选择性：只有成交的房源进了数据，挂了三个月没租出去的高价大户型不在里面，样本已经被结果筛过一遍。其三是控制变量选错了对象——把处理与结果的共同后果（例如``是否被中介重点推荐''）放进回归，反而会在原本无关的变量间制造出关联，这正是碰撞点效应。什么该控制、什么不该控制，判据在 \hyperref[sec:13-Machine-Learning/12-Causal-Inference/03]{``混杂、碰撞点与后门调整''}里；观察、干预与反事实三类问句的分层，见 \hyperref[sec:13-Machine-Learning/12-Causal-Inference/01]{``观察、干预与反事实回答不同问题''}与整个\hyperref[ch:120]{``因果推断：预测世界不等于改变世界''}一章。

另外两句常见的误读也值得挑明。``系数显著''说的是在零假设下观察到这么大的估计有多罕见，不说明效应大小有实际意义，样本量足够大时微不足道的差异也会显著。``\(R^2\) 高''说的是响应的变异被拟合值解释掉多少，与系数是否可信、与模型形式是否正确都没有直接关系；共线严重时 \(R^2\) 可以很高而每个系数的标准误都大到跨过零点。

\subsection{它至今没被换掉，靠的是可诊断}

线性模型的表达力早就被更复杂的模型盖过。它留在生产系统里的理由是另一件事：每个系数的来源可以指认，每个系数的不确定性可以量化，失效的方式已经被前人清点过，而且能在部署前查出来。

清点从固定预测目标开始。挂牌价、成交价与每平方米单价受不同程度的议价空间与季节波动影响，用挂牌价训练却拿成交价评价，会得到系统性的高估。切分之后，标准化所需的均值与标准差、类别变量的编码、缺失值的填补量，都只能从训练折算出来，否则测试误差里混进了它不该知道的信息（见 \hyperref[sec:13-Machine-Learning/01-Learning-Problems-and-Evaluation/04]{``数据划分、泄漏与分布偏移''}）。

拟合之后有四张图值得看。残差对拟合值作图，理想形状是无规律的水平带；呈漏斗形说明方差随水平变化，同方差假设失效，考虑加权最小二乘或对响应取对数；呈系统性弯曲说明线性形式漏掉了结构，该补特征或做变换，而不是去调优化器。杠杆值 \(h_{ii}\) 指出哪些样本的 \(x\) 远离数据中心，一套 500 平方米的房源在以 60 至 150 平方米为主的数据里会以不成比例的权重决定斜率，得单独判断它是录入错误还是部署时也会遇到的合法极端。条件数与 VIF 前面已经说过。训练误差与测试误差的落差则区分两类问题：落差大是记住了噪声，两者都高且接近是线性形式本身不够用。

最后是指标与代价对齐。均方误差按平方惩罚，500 元的误差记 250{,}000，100 元的误差记 10{,}000，前者是后者的 25 倍而误差本身只差 5 倍。如果平台对低价房源提供租金保障、高估低价房要现金赔付，那么两侧的真实损失是不对称的，训练时用的方便损失与部署后关心的损失必须分开来算（见\hyperref[sec:13-Machine-Learning/01-Learning-Problems-and-Evaluation/02]{``损失、风险、目标与指标不是同一个量''}）。

\subsection{线性指的是对参数线性}

顺带澄清一个术语。``线性模型''约束的是参数进入公式的方式，不要求对原始特征只能画直线。\(\widehat y=\beta_0+\beta_1x+\beta_2x^2+\beta_3x^3\) 仍是线性回归，因为待估参数以一次方式出现。换成交互项 \(\beta_4x_1x_2\)、样条基 \(B_1(x),\ldots,B_K(x)\)、径向基 \(\phi(\norm{x-c_k})\)，只要输出是参数与基函数值的点积，前面所有几何与分布结论都原样适用——你换掉的是基，没换掉线性结构。

麻烦出在外推。三次多项式能贴合训练范围内的弯曲走势，一旦 \(x\) 越出边界，曲线的行为就由 \(\beta_3x^3\) 主导，而高次项的系数恰好最依赖训练数据边缘那几个高杠杆点。更要紧的是，标准误覆盖的是``模型形式正确''前提下系数估计的随机误差，覆盖不了模型形式本身错了这件事。区间画得再窄，也只在假设成立的那个世界里成立。

下一节处理的正是本节留下的缺口：数据在某些方向上没有区分能力时，投影点稳而系数乱。正则化在那些方向上补一条数据之外的原则，告诉模型``分不清的时候你偏好什么样的解''，这一步买到的是方差和数值条件同时改善，卖掉的是无偏性。
